\paragraph{Qué modelo es mejor.} El ensamble obtiene el menor error (\mcensamble), seguido de LightGBM con los 33 rasgos (\mclgbmenriquecido) y de Ridge con los 33 rasgos (\mcridgeenriquecido). DeBERTa-v3 base (\mcdebertabase) y small (\mcdebertasmall) quedan por debajo de LightGBM, aunque claramente por encima de la línea base que solo usa la longitud (\mcbaselinelongitud). Todos los modelos superan a la línea base, de modo que el contenido del texto aporta información más allá del tamaño, pero la mejora sobre ella es moderada (de \mcbaselinelongitud\ a \mcensamble) porque la longitud explica buena parte de ambas notas.

\paragraph{Efecto de los rasgos nuevos.} Pasar de 18 a 33 rasgos (legibilidad, ortografía y similitud con la fuente) reduce el error de Ridge de \mcridgebase\ a \mcridgeenriquecido\ y el de LightGBM de \mclgbmbase\ a \mclgbmenriquecido. La mejora aparece en los cuatro textos fuente para ambos modelos, salvo Ridge en el texto \texttt{ebad26}, donde empeora ligeramente (tabla~\ref{tab:porprompt}), así que no es un efecto de un solo texto. Con LightGBM la ganancia viene sobre todo de \emph{wording}.

\paragraph{DeBERTa contra LightGBM.} DeBERTa predice \emph{content} casi tan bien como LightGBM (RMSE \rcdebertabase\ contra \rclgbmenriquecido), pero es peor en \emph{wording} (\rwdebertabase\ contra \rwlgbmenriquecido). Parte de esa diferencia es un sesgo: en los resúmenes fuera de muestra DeBERTa base sobreestima \emph{wording} en promedio $0{,}36$ puntos (LightGBM, $0{,}01$), y lo hace en los cuatro textos. Es plausible que, al entrenar con tres textos, aprenda un nivel de redacción propio de ellos que no se traslada al cuarto, pero con cuatro textos no podemos comprobar esa explicación. Su entrada no incluye el texto fuente, de modo que no puede medir directamente cuánto contenido del original recoge el resumen; los rasgos de similitud sí lo hacen. Además, el modelo grande no mejora al pequeño de forma consistente: es mejor en promedio, pero el orden se invierte en \texttt{39c16e}, y con 4 pliegues y una sola semilla esa diferencia no es concluyente. En costo, el ajuste de DeBERTa base tomó cerca de 500 segundos por pliegue en GPU contra menos de un segundo de LightGBM en CPU (figura~\ref{fig:tiempos}).

\paragraph{El ensamble.} El metamodelo (Ridge con coeficientes no negativos, evaluado también fuera de muestra) da a LightGBM el mayor peso (alrededor de $0{,}55$ para \emph{content} y $0{,}62$ para \emph{wording}), a DeBERTa base un peso intermedio ($0{,}28$ y $0{,}36$) y casi nada a Ridge. El ensamble mejora a LightGBM solo en $0{,}006$ (de \mclgbmenriquecido\ a \mcensamble), una diferencia pequeña frente a la variación entre textos (de $0{,}44$ a $0{,}62$ en el ensamble) y estimada con solo cuatro pliegues; por eso no la presentamos como una ventaja establecida. La ganancia por incorporar a DeBERTa es algo más clara si se compara con el ensamble que solo usa Ridge y LightGBM (\mcensamblerasgos), pero este último es peor que LightGBM solo, lo que indica que el apilado entre modelos parecidos no ayuda. Los errores de LightGBM y DeBERTa en \emph{wording} tienen correlación de $0{,}73$, es decir, se equivocan en buena medida en los mismos resúmenes, lo que limita lo que puede aportar la combinación.

\paragraph{Por texto fuente.} El texto más difícil para el ensamble es \texttt{814d6b} (\textit{The Third Wave}), con \mcensamble\ de promedio general pero $0{,}618$ en este texto. Ahí las notas de \emph{wording} son en promedio mucho más altas ($0{,}52$) que en los otros tres textos (entre $-0{,}30$ y $-0{,}07$); LightGBM subestima esa nota en $0{,}40$, mientras que DeBERTa la sobreestima solo en $0{,}07$, y por eso es el único texto donde DeBERTa base supera a LightGBM en \emph{wording}. Un cambio de nivel entre textos, causado quizá por la dificultad del texto o por el calificador, es el principal factor de error y no se corrige con rasgos del resumen.

\paragraph{Longitud y rasgos importantes.} El error crece con la longitud del resumen para todos los modelos (figura~\ref{fig:longitud}): en el quintil de resúmenes más cortos el ensamble tiene un error de $0{,}38$ y en el de los más largos de $0{,}68$. Los rasgos que más usa LightGBM para \emph{content} son el número de caracteres ($51\,\%$ de la ganancia) y de palabras ($24\,\%$); para \emph{wording} destacan también los caracteres ($27\,\%$), el solapamiento de bigramas ($17\,\%$) y el número de oraciones ($10\,\%$) (figura~\ref{fig:imp}). La importancia mide qué usa el modelo, no por qué una nota es alta; la longitud y las notas están correlacionadas, y los modelos aprovechan esa relación.
